Cieľ výstupu analýzy DNA získaných pomocou nanopórového sekvenovania je určiť vzťah medzi pôvodným DNA reťazcom a jeho pozorovanými čítaniami, ktoré vznikajú počas sekvenovania.
Namiesto priameho priraďovania konkrétneho referenčného reťazca k jednotlivým čítaniam je zvolený opačný postup: 
z dostupných nanopórových čítaní sa odhaduje najpravdepodobnejšia pôvodná sekvencia.
Tento odhad je založený na identifikácii dostatočne veľkej množiny podobných sekvencií, pri ktorých možno predpokladať spoločný pôvod.
Následne sa pre takúto skupinu vypočíta konsenzus, ktorý by mal byť čo najbližšie pôvodnému DNA reťazcu.

\textit{Konsenzus} je syntetizovaná sekvencia zložená z nukleotidov, ktoré sa v každej pozícii zarovnaných sekvencií vyskytujú s najväčšou frekvenciou.
Inými slovami, pre každú pozíciu v zarovnaní sa vyberie báza (A, C, G alebo T), ktorá sa v danej pozícii objavuje najčastejšie vo všetkých zarovnaných čítaniach.
Konsenzus tak reprezentuje najlepší odhad pôvodnej DNA sekvencie na základe dostupných pozorovaní a ich spoľahlivosti.

Keďže DNA reťazce sú dlhé a jednotlivé čítania obsahujú chyby, sekvencie sa môžu navzájom líšiť a byť vzájomne posunuté.
Z tohto dôvodu je potrebné vykonať viacnásobné zarovnanie sekvencií (multiple sequence alignment), aby sa konsenzus určoval z nukleotidov nachádzajúcich sa v rovnakých pozíciách.

Takto definovaný výstup umožňuje vytvoriť pravdepodobnostný model čítania DNA k-tíc (k-merov).
Tento model sa interpretuje ako komunikačný kanál a vychádza z matematickej teórie komunikácie \\cite{Shannon1948Communication}.
Namiesto jedného pevného výsledku sa každej k-tici priraďuje rozdelenie možných prečítaných sekvencií.
Formálne môžeme písať
$$
X=\Sigma^k,\quad Y=\Sigma^*,\quad K(y\mid x)=\Pr(Y=y\mid X=x),
$$
kde $\Sigma=\{A,C,G,T\}$, $x\in\Sigma^k$ je pôvodná k-tica a $y\in\Sigma^*$ je prečítaná sekvencia.
Tá môže mať aj inú dĺžku v dôsledku delécií alebo inzercií.

Pre konkrétnu k-ticu, napríklad $x=\mathrm{ATCG}$, môže byť rozdelenie nasledovné:
$$
\Pr(Y=\mathrm{ATCG}\mid X=x)=0.33,
$$
$$
\Pr(Y=\mathrm{ATTCG}\mid X=x)=0.20,
$$
$$
\Pr(Y=\mathrm{TCG}\mid X=x)=0.10,
$$
$$
\Pr(Y=\mathrm{AGCG}\mid X=x)=0.07.
$$
Zvyšná pravdepodobnosť patrí ostatným chybovým variantom tak, aby celkový súčet bol rovný 1.
Takýto model zachytáva substitúcie, inzercie a delécie a je kľúčový pre realistickú simuláciu chýb pri generovaní nových DNA sekvencií.

Zároveň platí, že jedna skupina podobných sekvencií spravidla neobsahuje dostatočný počet variácií na spoľahlivý odhad chýb pre všetky možné k-mery.
Počet k-tíc rastie exponenciálne podľa $4^k$; napríklad pre $k=4$ ide o $256$ možných 4-merov.
Pravdepodobnosť, že všetky tieto varianty budú dostatočne zastúpené v jednej skupine, je nízka, preto je potrebné objaviť viacero sekvenčných skupín.

Požadovaný výstup preto pozostáva z viacerých skupín navzájom zarovnaných DNA sekvencií a zodpovedajúcich konsenzuálnych sekvencií pre každú skupinu.
Takto definovaný výstup umožňuje dosiahnuť dostatočné pokrytie variácií čítaní odvodených z toho istého DNA reťazca.



